% Generated by member4_training_eval/make_ablation_text.py -- do not edit by hand.
\subsection{Ablation: Varying Strength and Modulation Separately}
\label{sec:ablation}

The main results compare a handful of configurations. The ablations below vary the two quantities of the method separately and over a grid, so that the conclusion does not rest on the choice of a single setting.  The $\bar{\alpha}$ column is the *measured* mean strength of each configuration, read back from the generator's own record of the blending field it applied, not the nominal parameter.

\paragraph{The strength sweep.} We vary the nominal strength $\alpha_0$ of the spatial variant while keeping the modulation depth at $\beta = 0.7$, and train the uniform AdaIN and the strength-matched uniform control at every level. The prediction of Section~\ref{sec:strength_energy} is that the accuracy follows the effective strength $\bar{\alpha}$ and that the three series lie on one curve.

\begin{center}
\scriptsize
\begin{tabular}{@{}llrrr@{}}
\toprule
$\alpha_0$ & configuration & $\bar{\alpha}$ & CrossAvg & $n$ \\
\midrule
$0.25$ & +Ours (spatial) & $0.187$ & $39.8\pm1.0$ & 5 \\
$0.25$ & +AdaIN (uniform) & $0.250$ & $42.0\pm1.0$ & 5 \\
$0.25$ & +AdaIN (matched) & $0.187$ & $40.9\pm0.9$ & 5 \\
$0.75$ & +Ours (spatial) & $0.560$ & $47.8\pm1.5$ & 5 \\
$0.75$ & +AdaIN (uniform) & $0.750$ & $51.1\pm2.3$ & 5 \\
$0.75$ & +AdaIN (matched) & $0.560$ & $47.7\pm1.4$ & 5 \\
$1.00$ & +Ours (spatial) & $0.747$ & $50.1\pm1.3$ & 5 \\
$1.00$ & +AdaIN (uniform) & $1.000$ & $54.2\pm1.8$ & 5 \\
$1.00$ & +AdaIN (matched) & $0.747$ & $50.5\pm1.6$ & 5 \\
\bottomrule
\end{tabular}
\end{center}

\paragraph{The modulation sweep.} Here the nominal strength is fixed at $\alpha_0 = 0.5$ and the modulation depth $\beta$ varies. At $\beta = 0$ the spatial variant degenerates to uniform AdaIN, so the sweep interpolates between the two extremes of the method. As $\beta$ grows, the modulation becomes stronger and the effective strength $\bar{\alpha}$ falls accordingly; the accuracy should follow the strength. The Edge column is the edge retention of the generated set, and the last column is the difference between the spatial variant and the matched control at the same effective strength.

\begin{center}
\scriptsize
\begin{tabular}{@{}llrrrr@{}}
\toprule
$\beta$ & configuration & $\bar{\alpha}$ & CrossAvg & $n$ & Edge \\
\midrule
$0.00$ & +Ours (spatial) & $0.500$ & $47.5\pm0.9$ & 5 & $0.3211$ \\
$0.00$ & +AdaIN (matched) & $0.500$ & $47.5\pm1.0$ & 5 & $0.3211$ \\
\bottomrule
\end{tabular}
\end{center}

\paragraph{The style-pool sweep.} The ten style references are split into two disjoint half-pools of five, and the spatial variant is trained against each. If the result depended on the particular paintings in the pool, the two half-pools would differ noticeably.
\label{sec:pool_sensitivity}

\begin{center}
\scriptsize
\begin{tabular}{@{}llrr@{}}
\toprule
Style pool & configuration & CrossAvg & $n$ \\
\midrule
full (10 styles) & +Ours (spatial) & $45.2\pm1.2$ & 5 \\
\midrule
pool A & +Ours (spatial) & $43.2\pm1.6$ & 5 \\
pool B & +Ours (spatial) & $43.7\pm1.1$ & 5 \\
\bottomrule
\end{tabular}
\end{center}

Figure~\ref{fig:pool_sensitivity} shows the same comparison with bootstrap intervals.  The two half-pools agree with each other to within the seed noise, and both fall slightly below the full pool: halving the pool removes style diversity, and the accuracy tracks that diversity rather than the modulation.  The comparison therefore does not rescue the spatial variant -- it moves every configuration in the same direction, which is what the strength account of Section~\ref{sec:strength_energy} predicts.

